\documentclass[10pt,a4paper]{amsart}
\usepackage[margin=19mm]{geometry}
\usepackage{amsmath,amssymb,mathtools}
\usepackage[scr=rsfs,cal=euler]{mathalfa}
\usepackage{xfrac}
\usepackage[unicode,hidelinks,bookmarks=false]{hyperref}
\newcommand{\ie}{i.e.\ }
\newcommand{\cf}{cf.\ }
\newcommand{\resp}{respectively\ }
\newcommand{\Subsection}{\subsection}
\providecommand{\nobreakdash}{\nobreak\hbox{-}}
\setlength{\emergencystretch}{2em}
\multlinegap0pt
\usepackage{lmodern}
\usepackage{mathtools}
\numberwithin{equation}{section}
\newtheorem{theorem}{Theorem}[section]
\newtheorem{lemma}[theorem]{Lemma}
\newtheorem{proposition}[theorem]{Proposition}
\newtheorem{remark}[theorem]{Remark}
\newtheorem{corollary}[theorem]{Corollary}
\newtheorem{definition}[theorem]{Definition}
\allowdisplaybreaks[1]
\setlength{\emergencystretch}{2em}
\newcommand{\bR}{\mathbb{R}}
\newcommand{\bC}{\mathbb{C}}
\newcommand{\bM}{\mathbb{M}}
\newcommand{\Mn}{\mathbb M_n}
\newcommand{\Mnp}{\mathbb M_n^+}
\newcommand{\Tr}{\operatorname{Tr}}
\newcommand{\diag}{\operatorname{diag}}
\newcommand{\spec}{\sigma}
\newcommand{\norm}[1]{\lVert #1\rVert}
\newcommand{\tnorm}[1]{\lVert #1\rVert_1}
\newcommand{\smod}[1]{\lvert #1\rvert_{\mathrm{sym}}}
\begin{document}
\title[Solutions to some open problems in matrix analysis]{Solutions to some open problems in matrix analysis}
\author{Mohamed Amine Aouichaoui; Eun-Young Lee}
\date{}
\maketitle
\noindent\emph{Source and licence.} Solutions to some open problems in matrix analysis. Version arXiv:2609.20094v2 (CC BY 4.0). This material is distributed under the Creative Commons Attribution 4.0 International licence, http://creativecommons.org/licenses/by/4.0/, as recorded in the arXiv OAI licence metadata for this exact version and shown on the abstract page https://arxiv.org/abs/2609.20094v2. Full rights notice: licenses/arxiv-2609.20094-CC-BY-4.0.txt.\par\smallskip
\noindent\textit{Editorial scope.} Counted unit: the original \texttt{theorem} environment \texttt{thm:radius-sharp} at source lines 1074--1087 together with its complete proof at lines 1089--1135; the delivered document keeps source lines 631--1136 so that every referenced label and every citation resolves inside it. Native editable source is the paper's own concatenated TeX; the AMS wrapper is editorial formatting only. No publisher class, upstream script or unrelated artwork is redistributed.
\begin{equation}\label{Rot}
\Tr f(A+B) \le \Tr f(A) + \Tr f(B).
\end{equation}
This trace inequality also holds for real-valued concave
functions $f$ on $[0,\infty)$ with $f(0)\geq0$.
The unitary-orbit inequality need not hold under these weaker
assumptions; see Zhang~\cite{TZhang26b}.

We use the following result of Fulton~\cite{Fulton} to prove
a stronger form of the Aujla--Bourin inequality.
Let $Z,A_1,\ldots,A_m\in\Mn$ be Hermitian.
There exist unitaries $U_1,\ldots,U_m\in\Mn$ such that
\[
 Z\leq\sum_{t=1}^m U_tA_tU_t^*
\]
if and only if
\[
 \sum_{k\in K}\lambda_k(Z)
 \leq
 \sum_{t=1}^m\sum_{i\in I_t}\lambda_i(A_t)
\]
for every Horn tuple $(I_1,\ldots,I_m;K)$ of subsets of
$\{1,\ldots,n\}$ with common cardinality $1\leq r\leq n$.
Here the tuple with
$I_1=\cdots=I_m=K=\{1,\ldots,n\}$ is included; it gives
\[
 \Tr Z\leq\sum_{t=1}^m\Tr A_t.
\]
For the definition of the Horn tuples, see~\cite[Section~2]{Fulton}.
We use this convention throughout the section.

We first prove the following result.

\begin{lemma}\label{lem:completion}
Let \(F,X,Y\in\Mnp\). Suppose that
\[
    \lambda_i(X)\leq\lambda_i(F),
    \qquad 1\leq i\leq n,
\]
and that
\[
    F\leq U_0 X U_0^*+V_0 Y V_0^*
\]
for some unitaries \(U_0,V_0\in\Mn\).
Then, there exist unitaries \(U,V\in\Mn\) such that
\begin{equation}\label{eq:completion}
    0\leq F-UXU^*\leq VYV^*.
\end{equation}
\end{lemma}

\begin{proof}

Put
\[
    \alpha=\lambda(X),\qquad
    \beta=\lambda(Y),\qquad
    \gamma=\lambda(F).
\]

We first prove that there exist positive semidefinite matrices
\(P,D\in\Mnp\) such that
\begin{equation}\label{eq:spectral-completion-revised}
    \lambda(P)=\alpha,\qquad
    \lambda_i(D)\leq\beta_i\quad(1\leq i\leq n),
    \qquad
    \lambda(P+D)=\gamma.
\end{equation}

We use Fulton's spectral characterization of the relation
\(C\leq A+B\); see~\cite[Theorem~1]{Fulton}.
By hypothesis,
\[
    F\leq U_0 X U_0^*+V_0 Y V_0^*.
\]
Hence the eigenvalue lists
\[
    \alpha=\lambda(X),\qquad
    \beta=\lambda(Y),\qquad
    \gamma=\lambda(F)
\]
form an admissible spectral triple for the relation \(C\leq A+B\).
By Fulton's spectral characterization, this is equivalent to saying
that \((\alpha,\beta,\gamma)\) satisfies all the majorized Horn
inequalities.



Moreover, by assumption,
\begin{equation}\label{eq:alpha-gamma}
    \alpha_i
    =\lambda_i(X)
    =\lambda_i(U_0 X U_0^*)
    \leq \lambda_i(F)
    =\gamma_i,
    \qquad 1\leq i\leq n.
\end{equation}

We first consider the case in which
\(\alpha,\beta,\gamma\) have rational entries. Multiplying all three
vectors by a common positive integer, and rescaling at the end, it is
enough to treat the case in which \(\alpha,\beta,\gamma\) are
partitions with at most \(n\) parts.

Write \(\mu\subseteq\nu\) if
\[
    \mu_i\leq\nu_i
    \qquad\text{for every }i,
\]
and denote by \(c_{\mu,\nu}^{\rho}\) the corresponding
Littlewood--Richardson coefficient.

Fulton's partition formulation of the majorized Horn theorem
\cite[Theorem~5 and the following remark]{Fulton} gives a partition
\(\widehat{\gamma}\) such that
\[
    \gamma\subseteq\widehat{\gamma}
    \qquad\text{and}\qquad
    c_{\alpha,\beta}^{\widehat{\gamma}}>0.
\]

Since $\alpha\subseteq\gamma\subseteq\widehat{\gamma}$,
Buch's refinement of Lemma~3, recorded in
\cite[footnote~2]{Fulton}, gives a partition
$\delta\subseteq\beta$ such that
\begin{equation}\label{eq:LR-refinement}
 c_{\alpha,\delta}^{\gamma}>0.
\end{equation}

The Horn--Littlewood--Richardson theorem now gives positive
semidefinite matrices \(P,D\in\Mnp\) such that
\[
    \lambda(P)=\alpha,\qquad
    \lambda(D)=\delta,\qquad
    \lambda(P+D)=\gamma.
\]

Since \(\delta\subseteq\beta\), we have
\[
    \lambda_i(D)=\delta_i\leq\beta_i,
    \qquad 1\leq i\leq n.
\]
Thus \eqref{eq:spectral-completion-revised} holds in the integral case,
and hence, after clearing denominators and rescaling, in the rational
case.

We now pass to arbitrary real spectra. Let \(\mathcal C\) denote the
set of triples \((a,b,c)\in\mathbb{R}^{3n}\) satisfying the majorized
Horn inequalities of Fulton, together with
\[
    a_1\geq\cdots\geq a_n\geq0,\qquad
    b_1\geq\cdots\geq b_n\geq0,\qquad
    c_1\geq\cdots\geq c_n\geq0,
\]
and
\[
    a_i\leq c_i,
    \qquad 1\leq i\leq n.
\]
All these inequalities have rational coefficients. Hence
\(\mathcal C\) is a rational polyhedral cone. 
By the preceding rational case, every rational point
\[
    (a,b,c)\in\mathcal C\cap\mathbb{Q}^{3n}
\]
is realized by positive semidefinite matrices \(P,D\) satisfying
\[
    \lambda(P)=a,\qquad
    \lambda_i(D)\leq b_i,\qquad
    \lambda(P+D)=c.
\]
Since \((\alpha,\beta,\gamma)\in\mathcal C\), and since rational
points are dense in the rational polyhedral cone \(\mathcal C\), there
exists a sequence of rational points
\[
    (\alpha^{(m)},\beta^{(m)},\gamma^{(m)})
    \in\mathcal C\cap\mathbb{Q}^{3n}
\]
such that
\[
    \alpha^{(m)}\longrightarrow\alpha,\qquad
    \beta^{(m)}\longrightarrow\beta,\qquad
    \gamma^{(m)}\longrightarrow\gamma.
\]
For each \(m\), choose positive semidefinite matrices
\(P_m,D_m\) such that
\[
    \lambda(P_m)=\alpha^{(m)},\qquad
    \lambda_i(D_m)\leq\beta_i^{(m)},\qquad
    \lambda(P_m+D_m)=\gamma^{(m)}.
\]
The sequences \((P_m)\) and \((D_m)\) are bounded. Indeed,
\[
    \|P_m\|
    =\lambda_1(P_m)
    =\alpha_1^{(m)}
\]
and
\[
    \|D_m\|
    =\lambda_1(D_m)
    \leq\beta_1^{(m)},
\]
and the right-hand sides remain bounded.
Since $\Mn$ is finite-dimensional and both sequences are bounded,
we may pass to a common subsequence such that
\[
    P_m\longrightarrow P,\qquad
    D_m\longrightarrow D
\]
for some \(P,D\in\Mnp\). 
Since the  eigenvalues depend continuously on the matrix, we
obtain
\[
    \lambda(P)
    =\lim_{m\to\infty}\lambda(P_m)
    =\alpha,
\]
and, for every \(i\),
\begin{equation}\label{eq:eigenvalue-continuity}
    \lambda_i(D)
    =\lim_{m\to\infty}\lambda_i(D_m)
    \leq
    \lim_{m\to\infty}\beta_i^{(m)}
    =\beta_i.
\end{equation}
Finally,
\[
    P_m+D_m\longrightarrow P+D,
\]
and therefore
\[
    \lambda(P+D)
    =
    \lim_{m\to\infty}\lambda(P_m+D_m)
    =
    \gamma.
\]
Thus \eqref{eq:spectral-completion-revised} holds for arbitrary real
spectra.

We now return to the original matrices \(F,X,Y\). Since
\(\lambda(P)=\lambda(X)\), there exists a unitary \(U\in\Mn\) such that
\[
    P=UXU^*.
\]
Since \(\lambda(P+D)=\lambda(F)\), there exists a unitary
\(W\in\Mn\) such that
\[
    F=W(P+D)W^*.
\]
Conjugating the matrices \(P,D\) by \(W\), and replacing \(U\) by
\(WU\), we may therefore assume, without loss of generality, that
\[
    F=P+D=UXU^*+D.
\]
In particular,
\[
    F-UXU^*=D\geq0.
\]
It remains to compare \(D\) with \(Y\). By \eqref{eq:eigenvalue-continuity},
\[
    \lambda_i(D)\leq\lambda_i(Y),
    \qquad 1\leq i\leq n.
\]
Equivalently, there exists a unitary \(V\in\Mn\) such that
\(
    D\leq VYV^*.
\)
Consequently,
\[
    0\leq F-UXU^*\leq VYV^*,
\]
which is \eqref{eq:completion}.
\end{proof}


\begin{theorem}\label{thm:double}

Let \(A,B\in\Mnp\) and let
\(f:[0,\infty)\to[0,\infty)\) be concave. Then, there exist unitaries
\(U,V\in\Mn\) such that
\begin{equation}\label{eq:double}
    0\leq f(A+B)-Uf(A)U^*
    \leq Vf(B)V^*.
\end{equation}
\end{theorem}

\begin{proof}
Set
\[
    F=f(A+B),\qquad
    X=f(A),\qquad
    Y=f(B).
\]
These matrices are positive semidefinite.

Since \(A\leq A+B\), we have
\[
    \lambda_i(A)\leq\lambda_i(A+B),
    \qquad 1\leq i\leq n.
\]
As a nonnegative concave function on \([0,\infty)\) is
nondecreasing, we infer
\[
    \lambda_i(X)
    =f(\lambda_i(A))
    \leq f(\lambda_i(A+B))
    =\lambda_i(F),
    \qquad 1\leq i\leq n.
\]
By the theorem of Aujla and Bourin
there exist unitaries
\(U_0,V_0\in\Mn\) such that
\[
    F\leq U_0 X U_0^*+V_0 Y V_0^*,
\]
and Lemma~\ref{lem:completion} now gives
\eqref{eq:double}.
\end{proof}

\begin{remark}
\label{rem:functional-horn}
By Fulton's spectral characterization, the Aujla--Bourin
inequality is equivalent to
\[
\sum_{k\in K} f\bigl(\lambda_k(A+B)\bigr)
\leq
\sum_{i\in I} f\bigl(\lambda_i(A)\bigr)
+
\sum_{j\in J} f\bigl(\lambda_j(B)\bigr)
\]
for every Horn tuple $(I,J;K)$.
Applying a concave function to the usual Horn inequalities does
not by itself give these inequalities. For results on concave
functions and Horn inequalities, see Zhou~\cite{JZhou26}.
\end{remark}



Theorem~\ref{thm:double} cannot in general be strengthened by requiring
$Vf(B)V^*\leq f(A+B)$ for the same unitaries $U,V$.
The following example shows this already in dimension $2$.

\begin{proposition}\label{prop:triple-counterexample}
Let
\[
 f(t)=\frac{t}{1+t}\quad(t\geq0),
 \qquad
 A=\begin{pmatrix}1&0\\0&0\end{pmatrix},
 \qquad
 B=\frac12\begin{pmatrix}1&1\\1&1\end{pmatrix}.
\]
There are no unitaries $U,V\in\mathbb M_2(\mathbb C)$ for which all three
inequalities
\begin{equation}\label{eq:triple-comparison}
 \begin{split}
 f(A+B)&\leq Uf(A)U^*+Vf(B)V^*,\\
 Uf(A)U^*&\leq f(A+B),\\
 Vf(B)V^*&\leq f(A+B)
 \end{split}
\end{equation}
hold. The matrices $A,B$ are real rank-one orthogonal projections,
and $f$ is nonnegative, smooth, strictly increasing and strictly
concave on $[0,\infty)$, with $f(0)=0$.
\end{proposition}

\begin{proof}
We have $A^2=A=A^*$, $B^2=B=B^*$ and
$\Tr A=\Tr B=1$. Also, the function $f(t)=1-\frac{1}{1+t}$ is obvioulsy strictly increasing and
strictly concave on $[0,\infty)$; it is also operator monotone
and operator concave.
Since $f(0)=0$ and $f(1)=1/2$,
\[
 f(A)=\frac12A,
 \qquad
 f(B)=\frac12B.
\]
Using $f(T)=I-(I+T)^{-1}$ for $T\geq0$, we obtain
\[
 F:=f(A+B)
 =I-(I+A+B)^{-1}
 =\frac17\begin{pmatrix}4&1\\1&2\end{pmatrix}>0,
 \qquad \Tr F=\frac67.
\]
Suppose that $U,V$ satisfy~\eqref{eq:triple-comparison}, and set
\[
 X=Uf(A)U^*,\qquad Y=Vf(B)V^*,
\]
and
\[
 R=F^{-1/2}XF^{-1/2},\qquad S=F^{-1/2}YF^{-1/2}.
\]
Both $R$ and $S$ have rank one. The three inequalities imply
\[
 0\leq R\leq I,\qquad 0\leq S\leq I,\qquad R+S\geq I.
\]
A positive semidefinite rank-one matrix bounded above by $I$ has
trace at most one. Consequently,
\[
 2\leq\Tr(R+S)=\Tr R+\Tr S\leq2.
\]
Thus $R+S-I\geq0$ has trace zero, so $R+S=I$.
Multiplying on both sides by $F^{1/2}$ gives $X+Y=F$.
Taking traces now yields the contradiction
\[
 1=\Tr f(A)+\Tr f(B)=\Tr(X+Y)=\Tr F=\frac67.
\]
\end{proof}

\section{The disk-radius coefficient for positive block matrices}
\label{sec:block}

Let
\[
\mathcal B=\begin{bmatrix}A&X\\X^*&B\end{bmatrix}\geq0,
\qquad A,B,X\in\Mn,
\]
and let $f:[0,\infty)\to[0,\infty)$ be concave.
For every Schatten $p$-norm, $1\leq p\leq\infty$, we have
\[
\|f(\mathcal B)\|_p\leq\|f(A)\|_p+\|f(B)\|_p;
\]
see~\cite[Corollary~3.5]{BL12}.
For $p=1$ and $X=A^{1/2}B^{1/2}$, this implies
Rotfel'd's inequality~\eqref{Rot}. 

\begingroup

In another direction,
\endgroup
Bourin and Lee proved the following estimate~\cite[Corollary~2.4]{BourinLee2022}. If
\[
 \mathcal A=\begin{bmatrix}A&N\\N^*&B\end{bmatrix}\geq0,
 \qquad A,B,N\in\Mn,
\]
$N$ is normal, and $\spec(N)$ lies in a disk of radius $r$, then
\begin{equation}\label{eq:block-bound}
 \lambda_{1+2j}(\mathcal A)
 \leq\lambda_{1+j}(A+B)+r,
 \qquad j=0,\ldots,n-1.
\end{equation}
Their Question~2.5 asks whether the coefficient of $r$ can be reduced and whether it is already best possible for $j=0$.


\begin{theorem}\label{thm:radius-sharp}
For every $r>0$ and $\varepsilon>0$, there are $A,B\in\mathbb M_3^+$ and a normal $N\in\mathbb M_3$ such that
\[
 \begin{bmatrix}A&N\\N^*&B\end{bmatrix}\geq0,
 \qquad
 \spec(N)\subseteq\{z\in\mathbb C:|z|\leq r\},
\]
and
\begin{equation}\label{eq:radius-sharp}
 \lambda_1\!\left(\begin{bmatrix}A&N\\N^*&B\end{bmatrix}\right)
 >\lambda_1(A+B)+r-\varepsilon.
\end{equation}
Hence the coefficient of $r$ in~\eqref{eq:block-bound} is best possible already for $j=0$.
\end{theorem}

\begin{proof}
Let $T>2$ and put
\[
 U=\begin{bmatrix}0&1&0\\0&0&1\\1&0&0\end{bmatrix},
 \qquad
 A_T=\diag\left(T,\frac2T,\frac T2\right),
 \qquad
 C_T=\diag\left(T,\frac T2,\frac2T\right).
\]
Define
\[
 B_T=U^*C_TU,
 \qquad
 \mathcal A_T=\begin{bmatrix}A_T&U\\U^*&B_T\end{bmatrix}.
\]
Then,
\[
 \diag(I_3,U)\,\mathcal A_T\,\diag(I_3,U^*)
 =\begin{bmatrix}A_T&I_3\\I_3&C_T\end{bmatrix}.
\]
After a permutation of the coordinates, the matrix on the right is the direct sum of
\[
 \begin{bmatrix}T&1\\1&T\end{bmatrix},
 \qquad
 \begin{bmatrix}2/T&1\\1&T/2\end{bmatrix},
 \qquad
 \begin{bmatrix}T/2&1\\1&2/T\end{bmatrix}.
\]
All three blocks are positive semidefinite. Their eigenvalues give
\[
 \spec(\mathcal A_T)
 =\left\{T+1,\ T-1,\ \frac T2+\frac2T,\ \frac T2+\frac2T,\ 0,\ 0\right\},
\]
so $\lambda_1(\mathcal A_T)=T+1$. On the other hand,
\[
 A_T+B_T=\diag\left(T+\frac2T,T+\frac2T,T\right).
\]
It follows that
\begin{equation}\label{eq:radius-gap}
 \lambda_1(\mathcal A_T)-\lambda_1(A_T+B_T)=1-\frac2T.
\end{equation}
Take $A=rA_T$, $B=rB_T$, and $N=rU$. Since $U$ is unitary, $N$ is normal and its spectrum lies in the closed disk of radius $r$ centered at zero. The block matrix is $r\mathcal A_T\geq0$. By~\eqref{eq:radius-gap}, the difference between the largest eigenvalues of $r\mathcal A_T$ and $A+B$ equals $r(1-2/T)$. Any
\[
 T>\max\{2,2r/\varepsilon\}
\]
therefore gives~\eqref{eq:radius-sharp}.
\end{proof}

\begin{thebibliography}{99}
\bibitem[BL12]{BL12} J.-C.~Bourin and E.-Y.~Lee, \emph{Unitary orbits of Hermitian operators with convex or concave functions}, Bull. Lond. Math. Soc. \textbf{44} (2012), no.~6, 1085--1102. \url{https://doi.org/10.1112/blms/bds080}
\bibitem[BourinLee2022]{BourinLee2022} J.-C. Bourin and E.-Y. Lee, \emph{Eigenvalue inequalities for positive block matrices with the inradius of the numerical range}, Internat. J. Math. \textbf{33} (2022), 2250009. \url{https://doi.org/10.1142/S0129167X22500094}.
\bibitem[Fulton]{Fulton} W. Fulton, \emph{Eigenvalues of majorized Hermitian matrices and Littlewood--Richardson coefficients}, Linear Algebra Appl. \textbf{319} (2000), 23--36. \url{https://doi.org/10.1016/S0024-3795(00)00218-4}.
\bibitem[JZhou26]{JZhou26} J.~Zhou, \emph{From majorization to Horn and unitary-orbit inequalities}, preprint, arXiv:2609.05854, 2026. \url{https://doi.org/10.48550/arXiv.2609.05854}.
\bibitem[TZhang26b]{TZhang26b} T.~Zhang, \emph{An operator triangle inequality for the quadratic symmetric modulus}, preprint, arXiv:2602.01463, 2026. \url{https://doi.org/10.48550/arXiv.2602.01463}
\end{thebibliography}
\end{document}
